\chapter{Conclusion}
\label{cha:conclusion}

This chapter summarizes the key findings of the thesis, discusses the limitations of the current approach, and outlines potential directions for future research in steerable terrain generation using diffusion models.

\section{Summary}
\label{sec:summary}

This thesis set out to bridge the gap between text-driven terrain generation and precise spatial control by integrating a ControlNet architecture into the MESA latent diffusion model. While text prompts offer high-level semantic control, they lack the geographic precision required for rigorous data-free geomodelling. To address this, specific geomorphological features, as described in Section~\ref{sec:terrain_feature_map_generation} were extracted from the DEM samples within the MajorTOM dataset. These feature maps were then used as conditioning hints. The implementation of a custom training pipeline demonstrated that the duplicated encoder blocks of the ControlNet could guide the generation process. The experiments confirmed that as the control weight increases, the spatial structure of the generated landscapes progressively conforms to the provided feature maps, successfully establishing a framework for steerable terrain synthesis.

The model achieved a mean IoU of $0.09$ for valley features, $0.066$ for ridge features, and $0.055$ for cliff features at a control weight of $1.4$. Comparing the first training run to the second with global normalization, valley feature mean IoU regressed from $0.117$ to $0.09$, ridge feature mIoU improved from $0.022$ to $0.066$, and cliff feature mIoU improved from $0.02$ to $0.055$. The increase in overall mIoU from $0.053$ to $0.075$ demonstrates the effectiveness of the global normalization strategy and longer training duration in enhancing the model's performance.

\section{Limitations and Future Work}
\label{sec:limitations}

Despite the successful integration of spatial conditioning, several limitations were identified during the evaluation. Even with the updated global normalization strategy, a mismatch remains between the DEM sample distribution and the base model's expected input distribution. Another issue is the noisy nature of mIoU scores for thin features like cliffs, which makes quantitative assessment of alignment with the conditioning hints difficult. A further limitation is the imbalance of control signals in the training data, which renders control of underrepresented features such as cliffs challenging. Possible avenues for future research include exploring improved normalization techniques to better align the data distributions and investigating alternative evaluation metrics that are more robust to thin features. Additionally, expanding the training dataset to include a more balanced representation of different features could improve controllability.

\section{Closing Remarks}
\label{sec:closing_remarks}

Overall, this thesis demonstrates the feasibility of integrating ControlNet into MESA for steerable terrain generation, while also highlighting key challenges and areas for improvement. By addressing these limitations, future work can further enhance the controllability and realism of diffusion-based terrain synthesis, paving the way for more powerful tools in data-free geomodelling.
